\documentclass[10pt,a4paper]{amsart}
\usepackage{amsmath,amssymb,amsthm,graphicx}
\usepackage{verbatim}
\newenvironment{reference}{\comment}{\endcomment}
%\newenvironment{reference}{}{}
\newenvironment{slogan}{\comment}{\endcomment}
\newenvironment{history}{\comment}{\endcomment}

% For commutative diagrams we use Xy-pic
\usepackage[all]{xy}

% We use 2cell for 2-commutative diagrams.
\xyoption{2cell}
% Theorem environments.
%
\theoremstyle{plain}
\newtheorem{theorem}[subsection]{Theorem}
\newtheorem{proposition}[subsection]{Proposition}
\newtheorem{lemma}[subsection]{Lemma}

\theoremstyle{definition}
\newtheorem{definition}[subsection]{Definition}
\newtheorem{example}[subsection]{Example}
\newtheorem{exercise}[subsection]{Exercise}
\newtheorem{situation}[subsection]{Situation}

\theoremstyle{remark}
\newtheorem{remark}[subsection]{Remark}
\newtheorem{remarks}[subsection]{Remarks}

\numberwithin{equation}{subsection}

% Macros
%
\def\lim{\mathop{\mathrm{lim}}\nolimits}
\def\colim{\mathop{\mathrm{colim}}\nolimits}
\def\Spec{\mathop{\mathrm{Spec}}}
\def\Hom{\mathop{\mathrm{Hom}}\nolimits}
\def\Ext{\mathop{\mathrm{Ext}}\nolimits}
\def\SheafHom{\mathop{\mathcal{H}\!\mathit{om}}\nolimits}
\def\SheafExt{\mathop{\mathcal{E}\!\mathit{xt}}\nolimits}
\def\Sch{\mathit{Sch}}
\def\Mor{\mathop{\mathrm{Mor}}\nolimits}
\def\Ob{\mathop{\mathrm{Ob}}\nolimits}
\def\Sh{\mathop{\mathit{Sh}}\nolimits}
\def\NL{\mathop{N\!L}\nolimits}
\def\CH{\mathop{\mathrm{CH}}\nolimits}
\def\proetale{{pro\text{-}\acute{e}tale}}
\def\etale{{\acute{e}tale}}
\def\QCoh{\mathit{QCoh}}
\def\Ker{\mathop{\mathrm{Ker}}}
\def\Im{\mathop{\mathrm{Im}}}
\def\Coker{\mathop{\mathrm{Coker}}}
\def\Coim{\mathop{\mathrm{Coim}}}

% Boxtimes
%
\DeclareMathSymbol{\boxtimes}{\mathbin}{AMSa}{"02}

%
% Macros for moduli stacks/spaces
%
\def\QCohstack{\mathcal{QC}\!\mathit{oh}}
\def\Cohstack{\mathcal{C}\!\mathit{oh}}
\def\Spacesstack{\mathcal{S}\!\mathit{paces}}
\def\Quotfunctor{\mathrm{Quot}}
\def\Hilbfunctor{\mathrm{Hilb}}
\def\Curvesstack{\mathcal{C}\!\mathit{urves}}
\def\Polarizedstack{\mathcal{P}\!\mathit{olarized}}
\def\Complexesstack{\mathcal{C}\!\mathit{omplexes}}
% \Pic is the operator that assigns to X its picard group, usage \Pic(X)
% \Picardstack_{X/B} denotes the Picard stack of X over B
% \Picardfunctor_{X/B} denotes the Picard functor of X over B
\def\Pic{\mathop{\mathrm{Pic}}\nolimits}
\def\Picardstack{\mathcal{P}\!\mathit{ic}}
\def\Picardfunctor{\mathrm{Pic}}
\def\Deformationcategory{\mathcal{D}\!\mathit{ef}}
\usepackage{cleveref}
\crefname{originallemma}{original lemma}{original lemmas}
\Crefname{originallemma}{Original lemma}{Original lemmas}
\makeatletter
\newlabel{categories-lemma-directed-category-system}{{0032}{source}{}{sourcepointer-1}{}}
\newlabel{categories-lemma-directed-category-system@cref}{{[originallemma][0][]0032}{[1][0][]source}{}{}{}}
\makeatother
\crefname{originallemma}{original lemma}{original lemmas}
\Crefname{originallemma}{Original lemma}{Original lemmas}
\makeatletter
\newlabel{lemma-module-colimit-fp}{{00HA}{source}{}{sourcepointer-2}{}}
\newlabel{lemma-module-colimit-fp@cref}{{[originallemma][0][]00HA}{[1][0][]source}{}{}{}}
\makeatother
\crefname{originallemma}{original lemma}{original lemmas}
\Crefname{originallemma}{Original lemma}{Original lemmas}
\makeatletter
\newlabel{lemma-ring-colimit-fp-category}{{0BUF}{source}{}{sourcepointer-3}{}}
\newlabel{lemma-ring-colimit-fp-category@cref}{{[originallemma][0][]0BUF}{[1][0][]source}{}{}{}}
\makeatother
\title{Every algebra is a directed colimit of finitely presented algebras, with surjective transitions in the finite type case}
\author{The Stacks Project Authors}
\date{stacks-project:algebra.tex:local-snapshot-20261012; GNU FDL 1.2 or later}
\begin{document}
\maketitle
\noindent\textit{Editorial note.} All delivered mathematical text is the author's own, copied from the pinned Stacks Project chapter file algebra.tex; no statement or proof step was rewritten, reordered or removed.\par\medskip
\noindent\textit{Editorial note.} Preamble provenance, as required for this source: the frozen snapshot stages the chapter's own preamble file beside it (preamble.tex, 217 lines, 8116 bytes, the same shared Stacks-wide preamble the published chow and etale-cohomology cards use); the first line of the pinned chapter file is \textbackslash{}input\{preamble\}, so no preamble text lives inside the chapter itself. The card takes that file's lines 8-18 (the verbatim package and the three \textbackslash{}comment-based environments reference/slogan/history at lines 9-12, section 13.2, together with the Xy-pic setup) and its lines 154 to the runtime end of the file (the native theorem declarations and mathematical macros). Not one character was changed, no package was added, no macro was defined, and the excerpt serves only as preamble, never as a span coordinate.\par\medskip
\noindent\textit{Editorial note.} The same-chapter pointers are located at their \textbackslash{}label lines inside the pinned chapter file; the cross-chapter pointers are located in the project's own tag list staged in the snapshot (sources/stacks-chow/tags.txt), at the verbatim lines carrying their canonical tags.\par\medskip
\noindent\textit{Editorial note.} The second proof inside the counted span is the explicit construction; the first proof is a citation of a categorical statement, and both stay verbatim inside the span.\par\medskip
\noindent\textit{Editorial note.} The counted proof refers to same-chapter statements by label; they are recorded as navigation stubs pointed at their \textbackslash{}label lines.\par\medskip
\noindent\textit{Allowed standard prerequisites (editorial).}\begin{itemize}
\item Finite presentation as finitely many generators and relations. Used in the explicit construction in the second proof (lines 32067-32094). Supplies the finitely presented approximation attached to a finite piece of data.
\item Finite type as a quotient of a polynomial algebra. Used when arranging the transition maps to be surjective (lines 32080-32093). Lets the index set be thinned so that the transition maps become surjections.
\end{itemize}\medskip
\ifdefined\hypertarget\hypertarget{sourcepointer-1}{}\fi
\noindent\textit{Source reference.} Original lemma 0032: sources/stacks-chow/tags.txt, source lines 119--119. Navigation pointer for \textbackslash{}ref\{categories-lemma-directed-category-system\} used inside the counted span (line 32065); the argument that consumes it is retained in the span.\par\medskip
\ifdefined\hypertarget\hypertarget{sourcepointer-2}{}\fi
\noindent\textit{Source reference.} Original lemma 00HA: sources/stacks-algebra/pinned-original-source.tex, source lines 1847--1847. Navigation pointer for \textbackslash{}ref\{lemma-module-colimit-fp\} used inside the counted span (line 32069); the argument that consumes it is retained in the span.\par\medskip
\ifdefined\hypertarget\hypertarget{sourcepointer-3}{}\fi
\noindent\textit{Source reference.} Original lemma 0BUF: sources/stacks-algebra/pinned-original-source.tex, source lines 31988--31988. Navigation pointer for \textbackslash{}ref\{lemma-ring-colimit-fp-category\} used inside the counted span (line 32064); the argument that consumes it is retained in the span.\par\medskip
\begin{lemma}
\label{lemma-ring-colimit-fp}
Let $R \to A$ be a ring map. There exists a directed system $A_\lambda$ of
$R$-algebras of finite presentation such that $A = \colim_\lambda A_\lambda$.
If $A$ is of finite type over $R$ we may arrange it so that all the
transition maps in the system of $A_\lambda$ are surjective.
\end{lemma}

\begin{proof}
The first proof is that this follows from
Lemma \ref{lemma-ring-colimit-fp-category} and
Categories, Lemma \ref{categories-lemma-directed-category-system}.

\medskip\noindent
Second proof.
Compare with the proof of Lemma \ref{lemma-module-colimit-fp}.
Consider any finite subset $S \subset A$, and any finite
collection of polynomial relations $E$ among the elements of $S$.
So each $s \in S$ corresponds to $x_s \in A$ and
each $e \in E$ consists of a polynomial
$f_e \in R[X_s; s\in S]$ such that $f_e(x_s) = 0$.
Let $A_{S, E} = R[X_s; s\in S]/(f_e; e\in E)$
which is a finitely presented $R$-algebra.
There are canonical maps $A_{S, E} \to A$.
If $S \subset S'$ and if the elements of
$E$ correspond, via the map $R[X_s; s \in S] \to R[X_s; s\in S']$,
to a subset of $E'$, then there is an obvious map
$A_{S, E} \to A_{S', E'}$ commuting with the
maps to $A$. Thus, setting $\Lambda$ equal the set of pairs
$(S, E)$ with ordering by inclusion as above, we get a
directed partially ordered set.
It is clear that the colimit of this directed system is $A$.

\medskip\noindent
For the last statement, suppose $A = R[x_1, \ldots, x_n]/I$.
In this case, consider the subset $\Lambda' \subset \Lambda$
consisting of those systems $(S, E)$ above
with $S = \{x_1, \ldots, x_n\}$. It is easy to see that
still $A = \colim_{\lambda' \in \Lambda'} A_{\lambda'}$.
Moreover, the transition maps are clearly surjective.
\end{proof}

\end{document}
